\section{The Myth of the Future}

\begin{quotex}
This is part of an article by \textbf{Guido de Giorgio}, titled The Instant and Eternity, first published in Diorama Filosofico in 1939.
\end{quotex}


We can say that the sacred is distinguished from the profane in what is essentially oriented toward the past to fix the stages of a procession which necessarily finds its culmination in a ``present''. This ``present'' is the metaphysical point where eternity throws itself, where the worlds dissolve in a fullness without limits, a duration without rhythm, a bliss without end. The present is eternity, the past is only the vestibule that leads toward, which inserts into eternity. To repeat, to retrace the whole cycle which is realized in the point means to carry with oneself the experience of the centuries, all cosmic evolution to unravel the framework in the pupil of God.

Faust could not stop the instant because he knew only the caducity of the instance, the immediate iridescence of illusion, the vertigo which submerges instead of transfiguring, the ``shadow of the flesh'', the labile and evanescent phantasm, not that which in God resides an infinite momentariness which is the mystery of the eternal now. Such are the two aspects of the ``instant'', according to which one places himself on the human or divine plane; it is about two apparently opposed and divergent points which mark two worlds, two rhythms, two realities, of which one is absolute, true, the other fallacious and illusory. Faust's words ``stop, you are so beautiful!'' is only a not very original lyrical substitute in the face of the unfathomed plenitude of the Ineffable where the mystery of the divine gestation takes place. The myth of purification through aesthetics is only the very fragile bridge thrown by modern imbecility onto the momentariness of the human-cosmic illusion in order to evade the positive certitude of the mystery, an impassable wall or else by the dizzying passing of the wing, i.e., of the Spirit of God.

That is why the modern world oscillates between a dead past and a nebulous future, between what is no longer and what will never be, except in the hope that anticipates and builds. Traditional wisdom, on the contrary, turns toward the past, lives it, enriches it, \emph{updates} it, inserting itself in it to lead it fully into the present and renew it in the \emph{ver aeternum} eternal spring which the Ancients attributed to the Golden Age, pointing to the perennial germination of the Truth, the multitude of transfiguring states, life which knows neither birth nor death, for it uncoils itself in the bliss of the realizing consciousness. But for moderns the past is past, dead, finished, concluded, closed, irremediable: ``le déjà vu, le déjà vécu'' already seen, already lived , says Bergson, according to a psychological orientation that clearly manifests all the nostalgic sentimentality of the small man horribly enslaved by his small world. So that between a dead past and a future not yet born, the twilight present swings, at once a cloudy waning and a too pale dawn, in sum a veritable pause in agony. And from this erroneous vision the \emph{myth of the future} is derived, the tension toward what is not, toward what will never be because in reality only the present, in absorbing the past, is the dynamic point, \emph{the whole bow of the ship} which faces the horizon \emph{but never reaches it}.


\flrightit{Posted on 2012-06-06 by Cologero}

